\documentclass{article}
\usepackage{amsmath,amssymb,hyperref}
\title{A modeled carrier departure qualifies a later shared radio receiver}
\date{26 November 2026}
\begin{document}
\maketitle
\paragraph{Result and scope.}
The earlier \href{https://github.com/eternal-november/eos-research-programmes/pull/624}{native dispatch replay}
assumed the admissible whole-history pairs of recipient and time.
This study obtains the useful timing exclusion mechanically from
finite movement primitives and an installed fixture rule binding
an authorization to an actual departure. In the positive case an
early report and an earlier owed report both actually reach their
simulator callbacks. The exclusion breaks if the carrier already
faces the exit, the gate is weakened to load set-down, or a clock
alignment admits the equality boundary of the radio ledger.
The warehouse-to-radio bridge remains authored; neither upstream
package implements this authority rule. This is a held-in
cross-layer claim under that declared rule, not field evidence or
full I1 or I2 qualification.

\paragraph{Frozen movement and authority source.}
Rebuild byte-pinned RWARE warehouse movement and shelf operations
using the manifest. An earlier loaded carrier begins at \((2,0)\);
exactly one native action is applied per quarter-second tick.
The custodian remains stopped for ticks 1 and 2 and performs a
native unload (toggle) at 3. The loaded position and shelf flags are
read from the native state, not constructed from issued commands.
The installed branch in the certificate either faces north
(requires a right turn, then one forward action to the named
east exit \((3,0)\)) or faces east (forward already possible).
A released shelf and an empty carrier at the pickup cell do not
prove departure. The fixture admits the installed continuation with
zero or one additional idle tick; each reaches the corresponding
exit within the due boundary. While the old carrier acts, a
second native warehouse carrier advances from \((1,3)\) to
\((2,3)\) at tick 4 on a separate path. A separate radio report
also receives a native callback at time zero.

On the original good route the owed report \(A\) retains its
separate radio recipient and its send is triggered by actual
set-down at tick 3. On the alternative degraded-route branch,
\(A\)'s only admitted recipient is the same physical radio
receiver used by report \(B\). The authored owner rule activates
this binding only when that exact old carrier has unloaded and
actually reached the east exit. The program obtains both event
times from native outputs and builds its support before B is
admitted. Replayed actual histories are separately rebuilt under
the chosen action sequence. It does not infer transport alignment,
old-route failure, or future authority from a model state alone.
The recipient/time correlation here is an \emph{enforced}
conditional law, and its application scope is assumed fixed for
the duration of the forecast. Any new permitted route, actor,
clock or effect that could affect this result needs requalification.

The lower bound is stronger than a timestamp collected after
execution. Before tick 3 there is no release in the declared
command schedule. If initially facing north at unload, no allowed
single action from \(\{\text{idle,left,right,forward}\}\) can
reach the east successor; only forward moves, and it cannot move
east in that orientation. A right turn and forward do achieve it.
If initially facing east, a single forward suffices.
Forty-two native prefix checks enumerate lengths zero, one and two
after unload across the two orientations. Longer words cannot
improve the first-exit lower bound. The positive service bound
additionally assumes the stated installed controller and bounded
stutter; the lower-bound test does not assume that arbitrary
faulty or idle executions ever depart.

\paragraph{Consumption and clock boundary.}
On the pinned SubT broker, a reliable early \(B\) of size \(b\)
consumes its shared sender and receiver ledgers in the native
radio accounting before its application callback. If \(A\) and \(B\) later share one of those
actual ledgers and the age is at most one second, \(A\)'s
size \(a\) fits precisely when \(a+b\le L\), for common
capacity \(L\). Exactly one second remains in the window;
after it the old debit is excluded. Individual sizes already
fit the cap. The calibrated fixture maps warehouse tick \(j\)
to radio-time \(j/4-\epsilon\), with \(\epsilon\) the relative
offset against the time-zero early action. The exact rational
test uses eighths and several fixed offsets, not a measured
clock synchronization guarantee. Certification before dispatch
can also use a qualified upper bound on \(\epsilon\) if the
rest of the horizon is stable.

For separate senders and any \(a+b>L\), early \(B\)
is sound on the long certain-gate support exactly when the
earliest possible alias occurs strictly after that second.
The no-stutter witness starts the binding at \(5/4-\epsilon\);
the permitted extra idle tick only delays it. Thus the
lower certificate allows a worst alignment
\(\epsilon_{\max}<1/4\), not equality. At \(\epsilon=1/4\)
the alias can enter the still-active row and only the planned
late slot is certified. A north-to-east movement right after
unload generates alias \(5/4\) with zero shift and two
three-byte reports fit a cap of five across those two times.
At the identical heading-independent unload tick, an east-facing
carrier would reach the exit at exactly one second, changing
the verdict. Shortcuts not listed in the authority/source
envelope would require evidence that they are impossible, not
a software assertion that the record is current.

For the direct comparison, use one permitted early \(B\) slot
at zero or a later slot \(t_A+1.25\).
The original \(A\) is owed by \(1.5\) seconds and B by \(2.75\)
seconds, so both schedules for B keep its deadline when its
send succeeds. A shared transmitter can still prevent early
B in the nonalias branch. All capacities,
payload effects and native callbacks are compared without mixing
different worlds or duplicating a ledger. The induced family
is not a physical mission with unknown timing or retry policies.

\paragraph{Comparisons and reuse boundary.}
The equally informed direct planner enumerates both B send
schedules using all source histories and native ledger predicates;
it ties the context rule. A conservative projection combining
separate recipient and time marginals retains no false approval
but inserts a false early-alias option. With the time-zero
north case, independent route at \(0.75\) and shared route
at \(1.25\) are each valid, but their cross pairing is not.
A snapshot arm sees only the current separate receiver.
A stale-proof control computes the certificate using departure
but actually installs a weaker unload-triggering gate: alias
can then occur at \(0.75-\epsilon\), and the old report
can really be lost after an early B. This is not a witness
against an intact actual departure gate. Using the weaker
record as a conservative bound while the stronger gate remains
installed only loses possible early service. Another control
uses the north certificate after changing the actual initial
heading, even though the gate definition is unchanged. The
shorter lawful path invalidates the old timing proof. This
identifies configuration and authority dependencies that must
remain pinned, not a technique for changing outside systems.

Across 4,608 cases per arm the exact joint version and the
safe marginal version lost no old callbacks, while the joint
version served B early in 152 more cases. The current-only
snapshot lost 323; stale gate and stale heading demonstrations
lost 76 and 38 respectively in their particular controls.
These are deterministic finite sweeps, not measurements of
failure frequency. Lean checks the projected motion and exact
clock boundary; complete source interpretation is separately
reviewed in the trace proof. No independent holdout or
architectural advantage over the equally informed comparator,
nor a deployment law connecting the two upstream plants.
\end{document}
